\subsection{模的系数}

设 E 是一个左 A-模。设 \(E^{*}\) 是 E 在 K 上的向量空间的对偶；我们赋予它它的自然右 A-模结构。给定 E 的元素 x 与 \(x^{*}\)（属于 \(E^{*}\)），我们用 \(c_{\mathrm{E}}(x, x^{*})\) 表示线性形式

\[
c _ {\mathrm{E}} (x, x ^ {*}): a \mapsto \langle x ^ {*}, a x \rangle\tag{5}
\]

（在 A 上）。这些线性形式称为 E 的系数。由 E 的系数生成的 \(\mathrm{A}^*\) 的子空间记作 \(\Theta_{\mathrm{E}}(\mathrm{A})\)。若 E 非零，则它非零。

若 F 是一个同构于 E 的 A-模，则 E 与 F 有相同的系数，且我们有 \(\Theta_{\mathrm{F}}(\mathrm{A}) = \Theta_{\mathrm{E}}(\mathrm{A})\)。把 \(E^{*}\) 视为一个左 \(A^{\circ}\) -模；E 的一个系数是 \(E^{*}\) 的一个系数，且 \(\Theta_{\mathrm{E}}(\mathrm{A}) \subset \Theta_{\mathrm{E}^{*}}(\mathrm{A}^{\circ})\)。因此，若 E 在 K 上有限维，则 E 与 \(E^{*}\) 有相同的系数，且我们有 \(\Theta_{\mathrm{E}}(\mathrm{A}) = \Theta_{\mathrm{E}^{*}}(\mathrm{A}^{\circ})\)。

设 \((\mathrm{M},\pi)\) 是代数 A 的一个线性表示。A-模 \(\mathrm{M}_{\pi}\) 的系数也称为表示 \(\pi\) 的系数。系数 \(c_{\mathrm{M}_{\pi}}(x,x^{*})\)（对 \(x\in \mathrm{M}\) 与 \(x^{*}\in \mathrm{M}^{*}\)）也记作 \(c_{\pi}(x,x^{*})\)；向量空间 \(\Theta_{\mathrm{M}_{\pi}}(\mathrm{A})\) 也记作 \(\Theta_{\pi}(\mathrm{A})\)。

注. —— 1) 设 M 在 K 上有限维。设 \((e_1, \ldots, e_n)\) 是 V 的一组基，\((e_1^*, \ldots, e_n^*)\) 是它的对偶基。用 \((\pi_{ij}(a))\) 表示 \(\pi(a)\) 关于 V 的基 \((e_1, \ldots, e_n)\) 的矩阵；我们有 \(\pi_{ij} = c_\pi(e_j, e_i^*)\)。映射 \(a \mapsto (\pi_{ij}(a))\) 是从 A 到 \(\mathbf{M}_n(\mathrm{K})\) 的一个同态；这样一个同态有时称为代数 A 的一个矩阵表示。

\begin{enumerate}
\def\labelenumi{\arabic{enumi})}
\setcounter{enumi}{1}
\tightlist
\item
  设 E 是一个 A-模，\(E'\) 是 E 的一个子模。由 II, §7, No.~5, 第 299 页，定理 5 的推论，我们有 \(\Theta_{\mathrm{E}'}(\mathrm{A}) \subset \Theta_{\mathrm{E}}(\mathrm{A})\) 与 \(\Theta_{\mathrm{E}/\mathrm{E}'}(\mathrm{A}) \subset \Theta_{\mathrm{E}}(\mathrm{A})\)。
\end{enumerate}

设 \(\gamma_{\mathrm{E}}\) 是从 \(\mathrm{E} \otimes_{\mathrm{K}} \mathrm{E}^{*}\) 到 \(\mathrm{A}^{*}\) 的、把 \(x \otimes x^{*}\) 映为 \(c_{\mathrm{E}}(x, x^{*})\) 的唯一的 K-线性映射。它是 (A,A)-双线性的，其像为 \(\Theta_{\mathrm{E}}(\mathrm{A})\)。我们由它导出 A-线性映射 \(c_{\mathrm{E}}^{\prime}: \mathrm{E} \to \operatorname{Hom}_{\mathrm{A}}(\mathrm{E}^{*}, \mathrm{A}^{*})\) 与 \(c_{\mathrm{E}}^{\prime \prime}: \mathrm{E}^{*} \to \operatorname{Hom}_{\mathrm{A}}(\mathrm{E}, \mathrm{A}^{*})\)，使得

\[
c _ {\mathrm{E}} ^ {\prime} (x) (x ^ {*}) = c _ {\mathrm{E}} (x, x ^ {*}) = c _ {\mathrm{E}} ^ {\prime \prime} (x ^ {*}) (x).
\]

用 \(\theta_{E}\) 表示 K-线性映射 \(\theta_{E}: E \otimes E^{*} \to \operatorname{End}_{K}(E)\)，它由关系

\[
\theta_ {\mathrm{E}} (x \otimes x ^ {*}) (y) = \langle x ^ {*}, y \rangle x
\]

（对 \(x, y \in E\) 与 \(x^{*} \in E^{*}\)）刻画。它的像是 E 的有限秩自同态所成的集 \(\operatorname{End}_{\mathrm{K}}^{\mathrm{f}}(\mathrm{E})\)（VIII, 第 463 页）。由迹的定义（出处同前），我们有

\[
\operatorname{Tr} \left(\theta_ {\mathrm{E}} \left(x \otimes x ^ {*}\right)\right) = \langle x ^ {*}, x \rangle
\]

（对 \(x \in E\) 与 \(x^{*} \in E^{*}\)）；因此我们有

\[
\langle \gamma_ {\mathrm{E}} (x \otimes x ^ {*}), a \rangle = \langle x ^ {*}, a x \rangle = \operatorname{Tr} (\theta_ {\mathrm{E}} (a x \otimes x ^ {*})) = \operatorname{Tr} (\theta_ {\mathrm{E}} (x \otimes x ^ {*}) a).
\]

这给出关系

\[
\langle \gamma_ {\mathrm{E}} (h), a \rangle = \operatorname{Tr} (\theta_ {\mathrm{E}} (h) \circ a _ {\mathrm{E}})
\]

（对 \(a \in \mathrm{A}\) 与 \(h \in \mathrm{E} \otimes_{\mathrm{K}} \mathrm{E}^*\)）。这证明 \(\Theta_{\mathrm{E}}(\mathrm{A})\) 是形如 \(a \mapsto \operatorname{Tr}(u \circ a_{\mathrm{E}})\) 的线性形式（定义在 \(\mathrm{A}\) 上）所成的集，其中 \(u\) 取遍 \(\operatorname{End}_{\mathrm{K}}^{\mathrm{f}}(\mathrm{E})\)。

引理 1. —— 映射 \(c_{\mathrm{E}}^{\prime \prime}\) 是双射。若 \(\mathrm{E}\) 有限维，则 \(c_{\mathrm{E}}^{\prime}\) 亦然。

当 F 是一个右 A-模且 G 是一个 K-向量空间时，我们在 II, §4, No.~1, 第 268 页，命题 1, a) 中定义了 K-向量空间的一个同构 \(\gamma\)，它从 \(\mathrm{Hom}_{\mathrm{K}}(\mathrm{F} \otimes_{\mathrm{A}} \mathrm{E}, \mathrm{G})\) 到 \(\mathrm{Hom}_{\mathrm{A}}(\mathrm{E}, \mathrm{Hom}_{\mathrm{K}}(\mathrm{F}, \mathrm{G}))\)。这个同构把 \(\varphi : \mathrm{F} \otimes_{\mathrm{A}} \mathrm{E} \to \mathrm{G}\) 映为同态 \(e \mapsto (f \mapsto \varphi(f \otimes e))\)。通过从 \(\mathrm{A}_d \otimes_{\mathrm{A}} \mathrm{E}\) 到 E 的典范同构（II, §3, No.~4, 第 249 页），\(\gamma\) 可等同于 \(c_{\mathrm{E}}''\)（当 \(\mathrm{F} = \mathrm{A}_d\) 且 \(\mathrm{G} = \mathrm{K}\) 时）。

类似地，II, §4, No.~1, 第 268 页，命题 1, b) 中定义的同构 \(\beta\) 特殊化为一个同构 \(\beta\)，它从 \(\mathrm{Hom}_{\mathrm{K}}(\mathrm{E}^{*}\otimes_{\mathrm{A}}\mathrm{A}_{s},\mathrm{K})\) 到 \(\mathrm{Hom}_{\mathrm{A}}(\mathrm{E}^{*},\mathrm{Hom}_{\mathrm{K}}(\mathrm{A}_{s},\mathrm{K}))\)。当 E 有限维时，E 可典范地等同于 \(\mathrm{Hom}_{\mathrm{K}}(\mathrm{E}^{*},\mathrm{K})\)，\(\mathrm{E}^{*}\) 可等同于 \(\mathrm{E}^{*}\otimes_{\mathrm{A}}\mathrm{A}_{s}\)；这时同构 \(\beta\) 等同于 \(c_{\mathrm{E}}^{\prime}\)。

命题 2. —— 设 E 是一个左 A-模。

\begin{enumerate}
\def\labelenumi{\alph{enumi})}
\item
  E 的系数所成的集是从 E 到 A* 的 A-线性映射的像的并。
\item
  此外，设 E 在 K 上的维数为 n。则左 A-模 \(\Theta_{\mathrm{E}}(\mathrm{A})\) 同构于 \(E^{n}\) 的一个商，A-模 E 同构于 \(\Theta_{\mathrm{E}}(\mathrm{A})^{n}\) 的一个子模，且 \(\Theta_{\mathrm{E}}(\mathrm{A})\) 的每个元素都是 \(E^{n}\) 的一个系数。
\end{enumerate}

断言 a) 由 \(c_{\mathrm{E}}^{\prime \prime}\) 的满射性得出。

我们来证明 b)。设 \((e_1^*,\ldots ,e_n^*)\) 是 \(\mathrm{E}^*\) 在 K 上的一组基。由于 \(\Theta_{\mathrm{E}}(\mathrm{A})\) 由 E 的系数生成，A-线性映射

\[
(x _ {1}, \ldots , x _ {n}) \mapsto \sum_ {i = 1} ^ {n} c _ {\mathrm{E}} (x _ {i}, e _ {i} ^ {*})
\]

从 \(\mathrm{E}^n\) 到 \(\Theta_{\mathrm{E}}(\mathrm{A})\) 是满射的。由 a)，\(\Theta_{\mathrm{E}}(\mathrm{A})\) 的每个元素都是 \(\mathrm{E}^n\) 的一个系数。此外，A-线性映射

\[
x \mapsto \left(c _ {\mathrm{E}} (x, e _ {1} ^ {*}), \ldots , c _ {\mathrm{E}} (x, e _ {n} ^ {*})\right)
\]

从 E 到 \(\Theta_{\mathrm{E}}(\mathrm{A})^{n}\) 是单射的；这就给出 b)。

注 3. —— 设 E 是 \(\mathrm{A}^*\) 的一个左 A-子模。设 \(\varepsilon\) 是 E 上的线性形式 \(y \mapsto y(1)\)。对 E 中每个 \(x\)，我们有 \(x = c_{\mathrm{E}}(x, \varepsilon)\)，故 E 是 \(\Theta_{\mathrm{E}}(\mathrm{A})\) 的一个 A-子模。

